\section{Tipo de investigación}

Con base en los objetivos, esta investigación se considera aplicada. La investigación aplicada, también conocida como práctica o empírica, se caracteriza porque busca poner en práctica o utilizar los conocimientos adquiridos. En este tipo, las consecuencias prácticas son un interés primordial del investigador \cite{aplicada}.

Además, se utiliza un diseño cuasiexperimental, porque se cambian variables independientes para ver cómo afectan los resultados de la simulación, sin tener un grupo control o asignar condiciones experimentales de forma aleatoria. 

Este tipo de investigación es útil porque ayuda a encontrar relaciones de causa y efecto entre los elementos del modelo y los comportamientos que surgen en el sistema, mientras se mantiene un entorno computacional que es controlado y se puede repetir.
